\section{\unclass Base External Simulation Interface}

\unclass This section is broken down into base level functions and  a section for the structs which those functions utilize.

\subsection{\unclass Functions}

\guidance{TIMELINE: This table should be drafted with all anticipated functions by SRR, and should be finalized by CDR. \\
It should include the intended functions accessed through the DLL. As development progresses, this should be updated.
}

\unclass This section contains a list of all base interface functions in the \gls{dll} with descriptions and explanations for each of them. These functions are common across all \gls{odessa} external simulation \gls{dll}s.


\tableOneNow{\unclasslabel}{\unclass Export Functions}{table1}{6.5in}{
    Function \\
    \hyperref[tab:Update]{bool Update()} \\
    \hyperref[tab:Cleanup]{bool Cleanup()} \\
    \hyperref[tab:DestroySystem]{void DestroySystem()} \\
    \hyperref[tab:GetCurrentTimestepStart]{double GetCurrentTimestepStart()} \\
    \hyperref[tab:GetCurrentTimestepEnd]{double GetCurrentTimestepEnd()} \\
    \hyperref[tab:GetCurrentTimestepInt]{int GetCurrentTimestepInt()} \\
    \hyperref[tab:GetSimTimeStepSize]{double GetSimTimeStepSize()} \\
    \hyperref[tab:AddExternalSystem]{void AddExternalSystem(AddExternalSystemInfo system\_info)} \\
    \hyperref[tab:AddExternalSystemExtended]{void AddExternalSystem(AddExternalSystemInfo\_Extended system\_info\_extended)} \\
    \hyperref[tab:UpdateExternalSystem]{void UpdateExternalSystem(ExternalSystemUpdateInfo system\_update)} \\
    \hyperref[tab:InitEnvironment]{void InitEnvironment()} \\
    \hyperref[tab:AddExternalSpdw]{void AddExternalSpdw(SpdwInit spdw\_init)} \\
    \hyperref[tab:GetVoacapAngles]{VoacapAngles GetVoacapAngles(source\_component\_id, dest\_component\_id, frequency\_hz)} \\
    \hyperref[tab:GetExternalSystemInfoExtended]{ExternalSystemInfoExtended GetExternalSystemInfo\_Extended()} \\
}
%
%\pagebreak\vspace{.35in}\vbox{
    \sectionthr{Update}

    \hspace{\parindent}\unclass The \textcolor{red5}{\textbf{Update}} function is a call applied to all \gls{odessa} objects within the \gls{dll}. First the \textit{System} updates itself, then the \textit{Environment}. \textcolor{red5}{\textbf{Update}} takes no input parameters and does not provide \gls{odessa} objects with any information about the rest of the simulation. The accompanying \textit{model specific} ICD should contain a \textcolor{red5}{\textbf{SetInput}} function that provides simulation data \textit{to} the \gls{odessa} object. In most cases the \textit{model specific} \textcolor{red5}{\textbf{SetInput}} function should be called before \textcolor{red5}{\textbf{Update}} is. This function throws a \textit{generic error} described below.
%    
    \hypertarget{genericerror}{\textit{Note: other functions reference this throw error, therefore it is written generically}.}
%
    The accompanying \textit{model specific} ICD should contain a function which initializes the \gls{odessa} system object in the \gls{dll}. If the \textcolor{red5}{\textbf{FUNCTION}} is called before the \textit{model specific} initialization function is called, \textcolor{red5}{\textbf{FUNCTION}} will throw a \textit{generic error}. This \textit{generic} throw will print an error message which reads: 
    
    ``Error: cannot call function \textcolor{red5}{\textbf{FUNCTION}} before system is initialized.''
%


    \functionTable{\unclasslabel}{Update}{Update}{bool Update()}
    {
             & No input value taken\\ 
    }{
        bool & A bool which always returns true, this provides no meaningful information\\
    }
%
%
%\pagebreak\vspace{.35in}\vbox{
    \sectionthr{Cleanup}

    \hspace{\parindent}\unclass The \textcolor{red5}{\textbf{Cleanup}} function is a call applied to all \gls{odessa} objects within the \gls{dll}. First the \textit{System} performs its end of timestep cleanup functionality, then any internal representations of external systems perform theirs. \textcolor{red5}{\textbf{Cleanup}} takes no input parameters and does not provide \gls{odessa} objects with any information about the rest of the simulation.\textcolor{red5}{\textbf{ Cleanup}} should be called at the end of every simulation timestep. This function throws a \hyperlink{genericerror}{\textit{generic error}} as described above.

    \functionTable{\unclasslabel}{Cleanup}{Cleanup}{bool Cleanup()}
    {
             & No input value taken\\
    }{
        bool & A bool which always returns true, this provides no meaningful information\\
    }

%
%\pagebreak\vspace{.35in}\vbox{
    \sectionthr{Destroy System}

    \hspace{\parindent}\unclass The \textcolor{red5}{\textbf{DestroySystem}} function is used at the very end of a simulation to release memory allocated within the \gls{dll}. This function does not throw errors.

    \functionTable{\unclasslabel}{Destroy System}{DestroySystem}{void DestroySystem()}
    {
             & No input value taken\\
    }{
        void & No return value given\\
    }

%
%\pagebreak\vspace{.35in}\vbox{
    \sectionthr{Get Current Timestep Start}

    \hspace{\parindent}\unclass The \textcolor{red5}{\textbf{GetCurrentTimestepStart}} returns the simulation time in seconds at the beginning of the current timestep. This function throws a \hyperlink{genericerror}{\textit{generic error}} as described above.

    \functionTable{\unclasslabel}{Get Current Timestep Start}{GetCurrentTimestepStart}{double GetCurrentTimestepStart()}
    {
               & No input value taken\\
    }{
        double & The beginning of the current simulation timestep in seconds\\
    }
%
%
%\pagebreak\vspace{.35in}\vbox{
    \sectionthr{Get Current Timestep End}

    \hspace{\parindent}\unclass The \textcolor{red5}{\textbf{GetCurrentTimestepEnd}} returns the simulation time in seconds at the end of the current timestep. This function throws a \hyperlink{genericerror}{\textit{generic error}} as described above.

    \functionTable{\unclasslabel}{Get Current Timestep End}{GetCurrentTimestepEnd}{double GetCurrentTimestepEnd()}
    {
               & No input value taken\\
    }{
        double & The end of the current simulation timestep in seconds\\
    }
%
%
%\pagebreak\vspace{.35in}\vbox{
    \sectionthr{Get Current Timestep Int}

    \hspace{\parindent}\unclass The \textcolor{red5}{\textbf{GetCurrentTimestepInt}} returns the current simulation timestep as an integer count. This function throws a \hyperlink{genericerror}{\textit{generic error}} as described above.

    \functionTable{\unclasslabel}{Get Current Timestep Int}{GetCurrentTimestepInt}{int GetCurrentTimestepInt()}
    {
               & No input value taken\\
    }{
        double & The current simulation timestep as an integer count\\
    }
%
%
%\pagebreak\vspace{.35in}\vbox{
    \sectionthr{Get Simulation Timestep Size}

    \hspace{\parindent}\unclass The \textcolor{red5}{\textbf{GetCurrentTimeStepSize}} function returns the length of the internal \gls{dll} timestep as a double in seconds. This value is set through a selected \textit{.scn} file and should remain the same throughout the simulation runtime. This function throws a \hyperlink{genericerror}{\textit{generic error}} as described above.

    \functionTable{\unclasslabel}{Get Current Timestep Size}{GetSimTimeStepSize}{double GetSimTimeStepSize()}
    {
               & No input value taken\\
    }{
        double & Returns the length of the internal \gls{dll} timestep in seconds\\
    }
%
%
%\pagebreak\vspace{.35in}\vbox{
    \sectionthr{Add External System}

    \hspace{\parindent}\unclass The \gls{dll} has an internal representation of entities from the external simulation of which the model or the environment need to be aware. These representations are added with the \textcolor{red5}{\textbf{AddExternalSystem}} function(s). This function applies inputs in values from the \textbf{\textit{\hyperref[tab:addtargetstruct]{Add External System Info Struct}}} to a new asset. This function should be called before the \textcolor{red5}{\textbf{InitEnvironment}} function so that the \gls{odessa} environment inside the \gls{dll} can properly register these representations. This function throws a \hyperlink{genericerror}{\textit{generic error}} as described above.

    \functionTable{\unclasslabel}{Add External System}{AddExternalSystem}{void AddExternalSystem(AddExternalSystemInfo system\_info)}
    {
        AddExternalSystemInfo system\_info & A \textbf{\textit{\hyperref[tab:addtargetstruct]{struct}}} containing relevant system info values\\
    }{
        void                               & No return value given\\
    }

%
%\pagebreak\vspace{.35in}\vbox{
    \sectionthr{Add External System (Extended)}

    \hspace{\parindent}\unclass The \gls{dll} has an internal representation of entities from the external simulation of which the model or the environment need to be aware. These representations are added with the \textcolor{red5}{\textbf{AddExternalSystem}} function(s). This function applies inputs in values from the \textbf{\textit{\hyperref[tab:addExternalSystemExtendedStruct]{Add External System Info Struct (Extended)}}} to a new asset. This function should be called before the \textcolor{red5}{\textbf{InitEnvironment}} function so that the \gls{odessa} environment inside the \gls{dll} can properly register these representations. This function throws a \hyperlink{genericerror}{\textit{generic error}} as described above.

    \functionTable{\unclasslabel}{Add External System (Extended)}{AddExternalSystemExtended}{void AddExternalSystem(AddExternalSystemInfo\_Extended system\_info\_extended)}
    {
        \scriptsize{AddExternalSystemInfo\_Extended} \normalsize{system\_info} & A \textbf{\textit{\hyperref[tab:addtargetstruct]{struct}}} containing relevant system info values\\
    }{
        void                                                                   & No return value given\\
    }
%
%
%\pagebreak\vspace{.35in}\vbox{
    \sectionthr{Update External System}

    \hspace{\parindent}\unclass The \textcolor{red5}{\textbf{UpdateExternalSystem}} function updates the positional data of an entity added to the simulation through the \gls{dll} with the \hyperref[tab:AddExternalSystem]{\textcolor{red5}{\textbf{AddExternalSystem}}} function(s). The external system is specified with the \textit{system\_id} member variable in the \textbf{\textit{\hyperref[tab:updatetargetstruct]{Update External System Info Struct}}}. This function throws a \hyperlink{genericerror}{\textit{generic error}} as described above.

    \functionTable{\unclasslabel}{Update External System}{UpdateExternalSystem}{void UpdateExternalSystem(ExternalSystemUpdateInfo system\_update)}
    {
        system\_update & A \textbf{\textit{\hyperref[tab:updatetargetstruct]{struct}}} containing relevant external system values\\
    }{
        void           & No return value given\\
    }

%
%\pagebreak\vspace{.35in}\vbox{
    \sectionthr{Initialize Environment}

    \hspace{\parindent}\unclass The \textcolor{red5}{\textbf{InitEnvironment}} function creates and initializes the \gls{odessa} environment within the \gls{dll} using input from a selected \textit{.scn} file. The \gls{odessa} environment is responsible for handling emitted \gls{spdw}s from external systems as they travel through the atmosphere. In order for \gls{spdw}s to be received by the model this \textcolor{red5}{\textbf{InitEnvironment}} function must be called before the \hyperref[tab:Update]{\textcolor{red5}{\textbf{Update}}} function. \hyperref[tab:AddExternalSystem]{\textcolor{red5}{\textbf{AddExternalSystem}}} calls should be performed before this function. This function throws a \hyperlink{genericerror}{\textit{generic error}} as described above.

    \functionTable{\unclasslabel}{Initialize Environment}{InitEnvironment}{void InitEnvironment()}
    {
             & No input value taken\\
    }{
        void & No return value given\\
    }
%
%
%\pagebreak\vspace{.35in}\vbox{
    \sectionthr{Add External SPDW}

    \hspace{\parindent}\unclass The \textcolor{red5}{\textbf{AddExternalSpdw}} injects an \gls{spdw} into the outbox of a generic transmitting component that represents the transmitter of an external system. This function should be called when a receiver needs to accept pulses from other models in the external simulation. In order to call it an external model must be added as an external system and the \textcolor{red5}{\textbf{InitEnvironment}} function must have been called; this will allow external models to appear in propagation paths.  This function throws a \hyperlink{genericerror}{\textit{generic error}} as described above.

    \functionTable{\unclasslabel}{Add External SPDW}{AddExternalSpdw}{void AddExternalSpdw(SpdwInit spdw\_init)}
    {
        spdw\_init & A \textbf{\textit{\hyperref[tab:spdwinitstruct]{struct}}} containing relevant \gls{spdw} values\\
    }{
        void & No return value given\\
    }

%
%
%\pagebreak\vspace{.35in}\vbox{
    \sectionthr{Get VOACAP Angles}

    \hspace{\parindent}\unclass The \textcolor{red5}{\textbf{GetVoacapAngles}} collects takeoff and landing angles calculated by \gls{voacap} for skywave propagation. See the UAM for additional information on \gls{voacap}. This function throws a \hyperlink{genericerror}{\textit{generic error}} as described above.

    \unclass \gls{odessa} system objects in the \gls{dll} will also throw an error if the \gls{voacap} input frequency is outside the operation limits. These limits are within two megahertz and thirty megahertz. This throw will print an error message which reads: 
    
    ``Error: VOACAP input frequency must be between 2 MHz and 30 MHz.''

    \functionTable{\unclasslabel}{Get VOACAP Angles}{GetVoacapAngles}{VoacapAngles GetVoacapAngles(source\_component\_id, dest\_component\_id, frequency\_hz)}
    {
        source\_component\_id & Integer which describes the component which the signal emits from\\
        dest\_compnent\_id & Integer which describes the component which the signal is to be received by\\
        frequency\_hz & Double which describes center frequency of an associated signal being propagated (hertz)\\
    }{
        VoacapAngles & A \textbf{\textit{\hyperref[tab:voacapanglesstruct]{struct}}} which contains the azimuthal and elevation angles for the takeoff and landing for skywave propagation\\
    }
%
%
%\pagebreak\vspace{.35in}\vbox{
    \sectionthr{Get External System Info Extended}

    \hspace{\parindent}\unclass The function \textcolor{red5}{\textbf{GetExternalSystemInfo\_Extended}} returns a struct representation of a model's system that can be given to other systems in the simulation. Other systems can use this struct to inform their copy of the environment code that a new system has entered the simulation. The environment can interact with this shell model system without having to know about the real model which lives in a different dll.

    \functionTable{\unclasslabel}{Get External System Info Extended}{GetExternalSystemInfoExtended}{ExternalSystemInfoExtended GetExternalSystemInfo\_Extended()}
    {
         & No input value taken \\
    }{
        ExternalSystemInfoExtended & A struct which contains the info needed to create a shell aproximation of this model in a different dll\\
    }
%
%
%
%
\pagebreak

\sectiontwo{Structs}

\unclass This section contains the structs that are common across all \gls{odessa} external simulation \gls{dll}s.
%
%\tableTwoNow{Add External System Information Struct}{2.5in, fg=TableBlue, l}{4in, l}{
\tableTwoNow{\unclasslabel}{\unclass External System Information Struct}{table1}{2.5in, fg=TableBlue, l}{4in, l}{
    \SetCell[c=2]{c}addTarget                                  & \\
    \SetCell[c=2]{altGra2, fg=black, font=\bfseries, c}Functions Using This Struct\\
    \SetCell[c=2]{c}FUNCTION NAMES                                  & \\
    \SetRow{altGra2, fg=black, font=\bfseries, c}Member & Description\\
    int system\_id                                              & Integer \gls{id} of the external system\\
    int entity\_id                                              & Integer \gls{id} of the entity within the external system\\
    int platform\_id                                            & Integer \gls{id} of the platform within the external system\\
    bool has\_tx                                                & Bool value describing whether the external system has a transmitter or not\\
    int interface\_id                                           & Integer \gls{id} of the interface within the external system. Only used if `has\_tx' is true\\
    int tx\_id                                                  & Integer \gls{id} of the transmitter within the external system. Only used if `has\_tx' is true\\
    int number\_of\_hits\_to\_death                             & Integer containing value of ``hits to kill'' an asset\\
    double rcs                                                  & Double containing the radar cross section, the apparent size of an asset in meters squared\\
    double initial\_speed                                       & Double containing initial speed in meters per second of an asset\\
    double initial\_heading                                     & Double containing initial directional facing of an asset in radians where 0.0 is due north\\
    std::array$<$double, 3$>$ dimensions                        & An array containing the dimensions of a target in length(m), width(m), and height(m) respectively\\
    std::array$<$double, 3$>$ initial\_lla                      & Initial \gls{lla} of a target in latitude(radians), longitude(radians), and altitude(m) respectively\\
    double operating\_frequency\_hz                             & Double containing the operating frequency in Hertz of the external system's transmitter. Only used if `has\_tx' is true\\
}
%
%\pagebreak

\tableTwoNow{\unclasslabel}{\unclass Add External System Information Struct (Extended)}{LABEL}{2.5in, fg=TableBlue, l}{4in, l}{
    \SetCell[c=2]{c}addExternalSystemExtended                                  & \\
	\SetCell[c=2]{altGra2, fg=black, font=\bfseries, c}Functions Using This Struct\\
	\SetCell[c=2]{c}FUNCTION NAMES                                  & \\
	\SetRow{altGra2, fg=black, font=\bfseries, c}Member & Description\\
    int system\_id                                              & Integer \gls{id} of the external system\\
    Color alignment                                             & Alignment of the external system (red, blue, or gray)\\
    std::vector$<$ExternalEntityInfo$>$ entity\_info            & Vector of external entity information structs to represent entities within the external system\\
}
%
\tableTwoNow{\unclasslabel}{\unclass External Entity Information Struct}{LABEL}{2.5in, fg=TableBlue, l}{4in, l}{
    \SetCell[c=2]{c}externalEntityInfo                                  & \\
	\SetCell[c=2]{altGra2, fg=black, font=\bfseries, c}Functions Using This Struct\\
	\SetCell[c=2]{c}FUNCTION NAMES                                  & \\
	\SetRow{altGra2, fg=black, font=\bfseries, c}Member & Description\\
    int entity\_id                                              & Integer \gls{id} of the entity within the external system.\\
    int platform\_id                                            & Integer \gls{id} of the platform within the external system.\\
    std::vector$<$ExternalAttachmentInfo$>$ attachment\_info    & Vector of external attachment info structs to represent the platform attachments of the external entity.\\
    int number\_of\_hits\_to\_death                             & Integer containing value of ``hits to kill'' an asset.\\
    double rcs                                                  & Double containing the radar cross section, the apparent size of an asset, in meters squared.\\
    double initial\_speed                                       & Double containing initial speed in meters per second of an asset.\\
    double initial\_heading                                     & Double containing initial directional facing of an asset in radians where 0.0 is due north.\\
    std::array$<$double, 3$>$ dimensions                        & An array containing the dimensions of a target in length(m), width(m), and height(m) respectively.\\
    std::array$<$double, 3$>$ initial\_lla                      & Initial \gls{lla} of a target in latitude(radians), longitude(radians), and altitude(m) respectively.\\
}
%
%\pagebreak
%
\tableTwoNow{\unclasslabel}{\unclass External Attachment Information Struct}{LABEL}{2.5in, fg=TableBlue, l}{4in, l}{
    \SetCell[c=2]{c}externalAttachmentInfo                                  & \\
	\SetCell[c=2]{altGra2, fg=black, font=\bfseries, c}Functions Using This Struct\\
	\SetCell[c=2]{c}FUNCTION NAMES                                  & \\
	\SetRow{altGra2, fg=black, font=\bfseries, c}Member & Description\\
    int interface\_id                                           & Integer \gls{id} of the external interface for this grouping of attachments\\
    std::vector$<$ExternalTransmitterInfo$>$ transmitter\_info  & Vector of external transmitter information structs to represent transmitters within this grouping of attachments\\
}

\tableTwoNow{\unclasslabel}{\unclass External Transmitter Information Struct}{LABEL}{2.5in, fg=TableBlue, l}{4in, l}{
    \SetCell[c=2]{c}externalTransmitterInfo                                  & \\
	\SetCell[c=2]{altGra2, fg=black, font=\bfseries, c}Functions Using This Struct\\
	\SetCell[c=2]{c}FUNCTION NAMES                                  & \\
	\SetRow{altGra2, fg=black, font=\bfseries, c}Member & Description\\
	    int transmitter\_id                                         & Integer \gls{id} of the transmitter\\
    double operating\_frequency\_hz                             & Double containing the operating frequency in Hertz of the transmitter\\
}
%
\tableTwoNow{\unclasslabel}{\unclass Update External System Information Struct}{LABEL}{2.5in, fg=TableBlue, l}{4in, l}{
    \SetCell[c=2]{c}ExternalSystemUpdateInfo                                  & \\
	\SetCell[c=2]{altGra2, fg=black, font=\bfseries, c}Functions Using This Struct\\
	\SetCell[c=2]{c}FUNCTION NAMES                                  & \\
	\SetRow{altGra2, fg=black, font=\bfseries, c}Member & Description\\
    int system\_id                                              & Integer \gls{id} of the external system\\
    double update\_time                                         & Simulation time in seconds for this update information\\
    std::vector$<$PlatformInfo$>$ new\_ecf\_states              & Updated target position and heading for each entity within the external system. These values are a described in the \textbf{\textit{\hyperref[tab:platforminfostruct]{Platform Information Struct}}}\\
}
%
%\pagebreak
%
\tableTwoNow{\unclasslabel}{\unclass Platform Information Struct}{LABEL}{2.5in, fg=TableBlue, l}{4in, l}{
    \SetCell[c=2]{c}PlatformInfo                                  & \\
	\SetCell[c=2]{altGra2, fg=black, font=\bfseries, c}Functions Using This Struct\\
	\SetCell[c=2]{c}FUNCTION NAMES                                  & \\
	\SetRow{altGra2, fg=black, font=\bfseries, c}Member & Description\\
    std::array$<$double, 3$>$ ecf\_position                     & Earth-centered, Earth-fixed coordinates of the platform.\linebreak Described as an array of doubles: X (meters), Y (meters), and Z (meters)\\
    std::array$<$double, 3$>$ ecf\_velocity                     & Earth-centered, Earth-fixed velocity of the platform.\linebreak Described as an array of doubles: X (meters per second), Y (meters per second), and Z (meters per second)\\
    std::array$<$double, 3$>$ ecf\_acceleration                 & Earth-centered, Earth-fixed acceleration of the platform.\linebreak Described as an array of doubles: X (meters per second squared), Y (meters per second squared), and Z (meters per second squared)\\
    std::array$<$double, 4$>$ pointing\_quaternion              & Current pointing quaternion of the platform.\linebreak Described as an array of doubles: q1 (unitless), q2 (unitless), q3 (unitless), and q4 (unitless)\\
}
%
\tableTwoNow{\unclasslabel}{\unclass SPDW Initialization Struct}{LABEL}{2.5in, fg=TableBlue, l}{4in, l}{
    \SetCell[c=2]{c}SpdwInit                                  & \\
	\SetCell[c=2]{altGra2, fg=black, font=\bfseries, c}Functions Using This Struct\\
	\SetCell[c=2]{c}FUNCTION NAMES                                  & \\
	\SetRow{altGra2, fg=black, font=\bfseries, c}Member & Description\\
    int external\_system\_id                                    & \gls{id} of the external system that generated the associated \gls{spdw}\\
    int sender\_id                                              & \gls{id} of the component within the external system that generated the associated \gls{spdw}\\
    int signal\_link                                            & An integer used to describe a class/group of pulses meant to be received as a burst\\
    int pulse\_number                                           & An integer used to describe the number of a pulse in a `signal\_link' group of pulses\\
    std::string metadata                                        & A string of key-value pairs describing model-specific metadata\\
    \gls{spdw}Info \gls{spdw}\_data                             & Qualities relevant to the specified \gls{spdw}, these values are a described in the \textbf{\textit{\hyperref[tab:spdwinfostruct]{SPDW Information Struct}}}\\
}
%
%\pagebreak
%
\tableTwoNow{\unclasslabel}{\unclass SPDW Information Struct}{LABEL}{2.5in, fg=TableBlue, l}{4in, l}{
    \SetCell[c=2]{c}SpdwInfo                                  & \\
	\SetCell[c=2]{altGra2, fg=black, font=\bfseries, c}Functions Using This Struct\\
	\SetCell[c=2]{c}FUNCTION NAMES                                  & \\
	\SetRow{altGra2, fg=black, font=\bfseries, c}Member & Description\\
    double transmit\_time\_secs                                 & Double which describes the transmit time of an signal (seconds)\\
    double pri\_secs                                            & Double which describes the Pulse Repetition Interval of an associated signal (seconds)\\
    double pulse\_width\_secs                                   & Double which describes pulse width of an associated signal (seconds)\\
    double erp\_watts                                           & Double which describes effective radiate power of an associated signal. This value is calculated post-environmental effects (watts)\\
    double center\_freq\_hertz                                  & Double which describes operational center frequency of an associated signal (hertz)\\
    double signal\_bandwidth\_hertz                             & Double which describes the operational bandwidth of an associated signal (hertz)\\
    double transmit\_phase\_rads                                & Double which describes the phase of an output antenna skin of an associated signal (radians)\\
    CoherencyType coherency                                     & Enum which describes the coherency of an associated signal. Includes coherent and noncoherent\\
    SkirtShape skirt\_shape                                     & Enum which describes the skirt shape of an associated signal\\
    JammingTypes jamming\_type                                  & Enum which describes the jamming type of an associated signal \textit{Note: for jamming \gls{spdw}s only}\\
    Polarization polarization                                   & Enum which describes the polarization of the transmitting antenna of an associated signal\\
    PulseModulation modulation\_on\_pulse                       & Enum which describes the modulation schema on the pulse of an associated signal\\
    PulseModulationSpec modulation\_spec                        & Enum which describes advanced parameters of the modulation of an associated signal\\
    TransmitMode transmit\_mode                                 & Enum which describes the transmit mode of an associated signal\\
}
%
%\pagebreak
%
\tableTwoNow{\unclasslabel}{\unclass VOACAP Angles Struct}{LABEL}{2.5in, fg=TableBlue, l}{4in, l}{
    \SetCell[c=2]{c}VoacapAngles                                  & \\
	\SetCell[c=2]{altGra2, fg=black, font=\bfseries, c}Functions Using This Struct\\
	\SetCell[c=2]{c}FUNCTION NAMES                                  & \\
	\SetRow{altGra2, fg=black, font=\bfseries, c}Member & Description\\
    AzimuthElevation landing\_rads                              & Landing angle from source component to destination component as calculated by \gls{voacap} (radians)\linebreak It is composed of by the \textbf{\textit{\hyperref[tab:azimuthelevationstruct]{Azimuth and Elevation Struct}}}\\
    AzimuthElevation takeoff\_rads                              & Takeoff angle from source component to destination component as calculated by \gls{voacap} (radians)\linebreak It is composed of by the \textbf{\textit{\hyperref[tab:azimuthelevationstruct]{Azimuth and Elevation Struct}}}\\
}
%
\tableTwoNow{\unclasslabel}{\unclass Azimuth and Elevation Struct}{LABEL}{2.5in, fg=TableBlue, l}{4in, l}{
    \SetCell[c=2]{c}AzimuthElevation                                  & \\
	\SetCell[c=2]{altGra2, fg=black, font=\bfseries, c}Functions Using This Struct\\
	\SetCell[c=2]{c}FUNCTION NAMES                                  & \\
	\SetRow{altGra2, fg=black, font=\bfseries, c}Member & Description\\
    double azimuth                                              & Double describing the azimuth angle (radians)\\
    double elevation                                            & Double describing the elevation angle (radians)\\
}
%
%%% - - -   UNCLASSIFIED   - - - %%%